\section{Introduction}
\label{sec:intro}

Students usually meet quantum error correction as a list of facts: a code is a subspace, a stabilizer is a Pauli operator, a syndrome is a bit string, and a decoder turns the bit string into a correction. The facts are correct, but they hide the part that makes fault tolerance hard in practice, namely that every protective measure is itself built out of noisy gates and has to pay for itself. The question posed by the Qiskit Fall Fest 2026 challenge~\cite{qff2026prompt} is how to let a student feel that tension directly. Our answer is a game in which every move is a design decision about a real error-detection scheme and every score is the outcome of a Qiskit simulation of that scheme.

The scheme is low-overhead error detection with spacetime codes~\cite{martiel2025spacetime}, available in Qiskit as the \texttt{qiskit-paulice} add-on~\cite{paulice,ibm_tutorial_spacetime}. A spacetime check is a single ancilla qubit that touches one target qubit of a Clifford circuit at several moments, applying a controlled Pauli each time, and is then measured. Combined with the parity of some output bits, the ancilla outcome is deterministic on an ideal run. A non-trivial outcome therefore flags an error, and the shot is discarded. The construction descends from coherent Pauli checks~\cite{vandenberg2023cpc} and coherent parity check codes~\cite{roffe2018cpc,chancellor2023graphical,eczoo_cpc}, but it spreads each check over time instead of measuring a heavy operator at one instant, which keeps it compatible with sparse hardware connectivity. IBM has used the same idea to sample from a 64-qubit, depth-73 Clifford circuit doped with 314 $T$ gates while certifying a lower bound on the fidelity~\cite{martiel2026sampling,ibm_blog_trusted}.

Where does such a scheme sit? Following the distinction drawn in the challenge text and in Ref.~\cite{ibm_blog_emc}, \emph{error detection} flags the errors its checks can see and throws those shots away. It does not repair the state, and errors that the checks cannot see go through untouched. \emph{Error mitigation} leaves the hardware as it is and corrects the estimate of an observable in post-processing, at the cost of extra circuits and shots~\cite{temme2017error}. \emph{Error correction} encodes logical qubits in many physical qubits, measures stabilizers repeatedly and uses a decoder to undo errors while the computation runs~\cite{shor1995scheme,fowler2012surface,krinner2022realizing,google2025below}. Our game lives entirely in the first category. The player never corrects anything; the only currency is the fraction of shots that survive post-selection.

\paragraph*{Design goal.} The challenge asks for a game whose rules correspond to the decisions of the error-detection workflow and whose score comes from computed numbers, not from hand-set values. We add two requirements of our own. First, the game should contain at least one trade-off that a strong student would not guess before playing. Second, the numbers behind the game should be good enough to stand as a small research study: exact where possible, validated against shot-based sampling, and averaged over several random circuits.

\paragraph*{Contributions.}
\begin{enumerate}[leftmargin=*,itemsep=2pt]
  \item \textbf{A game whose rules are the physics.} \emph{Spacetime Sentinels} asks the player to spend an ancilla budget on spacetime checks, to inject a quota of $T$ gates, and to choose how the $T$ gates are placed relative to the checks. The rule-to-physics map is given in Table~\ref{tab:map}. The game runs in the browser and in the terminal.
  \item \textbf{Every score from Qiskit.} Checks are found with \texttt{add\_pauli\_checks}, ideal states come from stabilizer or statevector simulation, and noisy post-selected states come from exact Qiskit Aer density-matrix simulation. We evaluated all \NumMoves{} legal moves of the game under two noise models, one taken from the challenge text and one derived from the calibration data of a Heron backend.
  \item \textbf{A principled score.} We combine post-selected fidelity and kept fraction into a fidelity rate $R_\beta=F_{\mathrm{ps}}\,\kappa^{\beta}$. For $\beta=1/2$ this is the signal-to-noise ratio per square root of a raw shot, and $R_{1/2}^{-2}$ is the sampling overhead of detection followed by rescaling.
  \item \textbf{Three non-obvious findings, backed by numbers.} (i) Checks do not pay for themselves on shallow circuits; there is a break-even depth. (ii) $T$ gates are free on the hardware but expensive for the code: every check removes the spacetime locations where a $T$ gate can be placed without breaking a check, so protection and magic compete for the same wires. (iii) A check broken by a $T$ gate is worse than no check at all, and it keeps hurting even after the player stops listening to it.
\end{enumerate}

\paragraph*{Related games.} Litinski's \emph{A Game of Surface Codes}~\cite{litinski2019game,pennylane_game} turns lattice surgery into a board game with explicit space and time costs, and was the template for our spacetime board. \emph{Quantum Poker}~\cite{fuchs2019poker} uses a game to benchmark error mitigation, Decodoku and its genetic-algorithm solver~\cite{wootton2016minimal,genetic_decodoku} turn decoding into a puzzle, and Ref.~\cite{guerrero2025nash} uses game theory to search for codes. To our knowledge no existing game is built on error detection with spacetime codes, nor on the interplay between check placement and non-Clifford doping.
